%% ma: L11-B16-0044
%% lop: 11
%% chuong: 5
%% bai: 16
%% dang: 11.16.1
%% loai: TN4
%% muc: VD
%% dap_an: "A"
%% nguon: "tai-lieu-hoc-tap-mon-toan-lop-11-theo-chuong-trinh-giao-duc-pho-thong-2018.pdf (Lớp Toán Thầy Duy, Phạm Lê Duy, THPT Chu Văn An), Chương 3 Bài 2, Phần 1, Câu 32 (tr.369)"
%% kien_thuc: "Lớp 11 Bài 16 (giới hạn của hàm số)"
%% trang_thai: cho-duyet
%% ly_do: "Dạng mới đề xuất (11.16.x), chờ giáo viên duyệt dạng; tài liệu gốc không có lời giải mẫu, đáp án do kho tự giải và kiểm bằng sympy."
\begin{ex}
Giá trị của $\lim\limits_{x\to1}\dfrac{x^{2018}+x-2}{x^{2017}+x-2}$ bằng $\dfrac ab$, với $\dfrac ab$ là phân số tối giản. Tính giá trị của $a^2-b^2$.
\choice{\True $4037$}{$4035$}{$-4035$}{$4033$}
\loigiai{\sloppy Tử $=(x^{2018}-1)+(x-1)=(x-1)\left(x^{2017}+x^{2016}+\cdots+1+1\right)$ có tổng $2018$ số hạng $x^k$ cộng thêm $1$; tại $x=1$ thừa số còn lại bằng $2018+1=2019$. Tương tự mẫu: $2017+1=2018$. Giới hạn $=\dfrac{2019}{2018}$ (tối giản), nên $a^2-b^2=(a-b)(a+b)=1\cdot4037=4037$.}
\end{ex}
